\documentclass[11pt]{article}
\usepackage[margin=0.9in]{geometry}
\usepackage[T1]{fontenc}
\usepackage[utf8]{inputenc}
\usepackage{amsmath,amssymb}
\usepackage{enumitem}
\setlist[itemize]{topsep=3pt,itemsep=2pt,parsep=0pt}
\setlist[enumerate]{topsep=3pt,itemsep=2pt,parsep=0pt}

\title{Stage 00 Computational Project Map v1}
\author{Reaction-boundary volatility surface project}
\date{5 July 2026}

\begin{document}
\maketitle

\section*{Artefact Status}
Draft computational map. Stage 1 and Stage 1A only. No implementation,
figures, tables, numerical verification, or mathematical verification has
been performed.

\section*{Source Theory}
Primary source:
\begin{itemize}
  \item \texttt{CATG-BSMFromAdjointProjection-FMF-v1.8(3).tex}
  \item \texttt{BSMFromAdjointProjection-FMF-v1.8(3).bib}
\end{itemize}

\section*{Project Aim}
Develop a Python-first computational workflow that simulates a reasonable
approximation to the reaction-boundary volatility implied by the latent
order-book theory and generates a volatility surface from the main-result
closure. This is a simulation-only project at this stage.

\section*{Scope Decision}
Routing: simulation only.

No empirical data, option quotes, calibration routines, data cleaning,
sampling-window choices, or fitted market surfaces are part of the current
workflow. A later human comparison to a calibrated surface may be useful,
but it is outside this project unless explicitly added by a new control
decision.

\section*{Source Sections To Use}
\begin{itemize}
  \item Main-result operational closure, identified by the user as the
  target Eq.~18 rather than the earlier spectral/cumulant formula.
  \item Clock projection for a deterministic activity clock.
  \item Non-unique time and adjoint-real projections, only as a later
  caveat layer unless the simulation explicitly models clock ambiguity.
  \item Reaction-boundary response and Green-function cumulant as source
  motivation, not as first-pass direct simulation targets.
  \item Appendix finite-cutoff spectral form only as a deferred numerical
  refinement.
\end{itemize}

\section*{Main Computational Objects}
\begin{itemize}
  \item Primary target: the main-result operational reaction-boundary
  variance closure, using the source label \texttt{eq:operational-result}
  or its equivalent final numbered equation:
  \[
    a_u^{(\Delta)}(S,u)
    \simeq
    \left[\frac{\mathcal A_{\rm eff}(S,u)}{\mathcal L_u(S,u)^2}\right]
    \Delta^{-\gamma(S,u)}
    \mathcal F_{\gamma(S,u)}(\nu_u(S,u)\Delta).
  \]
  \item Deterministic-clock local variance projection
  \[
    \sigma_{\rm loc}^2(S,t)=\alpha(t)\Xi(S,U(t);\Delta).
  \]
  \item The same object written as
  \(\Xi(S,u;\Delta)\) for deterministic-clock projection.
  \item Volatility surface representation over stock price or log-price
  state and maturity after choosing a surface-generation convention.
\end{itemize}

\section*{Key Modelling Questions For Stage 2--5}
\begin{itemize}
  \item Response function: do not simulate the full response path first.
  First choose a controlled approximation or numerical evaluation of
  \(\mathcal F_\gamma(z)\), with \(z=\nu\Delta\).
  The two asymptotic cases in the response-function equation identified
  by the user as Eq.~22 should both be simulated: the low-resilience /
  small-\(z\) case and the high-resilience / large-\(z\) case.
  \item Stock-price dependence: the first plausible source is
  state-dependent liquidity \(\mathcal L(S)\), since the kernel contains
  \(1/\mathcal L(S)^2\). State-dependent \(\mathcal A_{\rm eff}(S)\),
  \(\nu(S)\), \(\gamma(S)\), or \(\alpha(t)\) may be introduced only as
  explicit modelling choices.
  \item Surface axes: decide whether the first surface uses \(S=\log x\),
  stock price \(x\), or log-moneyness. The project should not hide this
  convention.
\end{itemize}

\section*{Prudent Initial Outputs}
\begin{enumerate}
  \item A standalone parameter-range table.
  \item A two-panel figure: surface plot of simulated local volatility and
  contour projection of the same surface, with the two Eq.~22 response
  regimes either shown as separate panels or generated as matched outputs.
  \item A response-regime diagnostic plot or table comparing the two
  asymptotic cases of \(\mathcal F_\gamma(z)\).
  \item Optional liquidity-profile plot, only if state-dependent
  \(\mathcal L(S)\) is used to create stock-price dependence.
\end{enumerate}

\section*{Claims Boundary}
The computation should be labelled as simulation or illustration.
Computational evidence must not be treated as proof of the theory, and
visual surface shape must not be described as empirical validation.

\end{document}
